\Needspace{14\baselineskip}
\SourceLeaf{102}%
\begin{center}
BEMERKUNGEN ZU DEN UNTERSUCHUNGEN ÜBER DIE AUS DEM INTEGRAL
$\displaystyle\int\frac{dx}{\BareRoot{(1+x^3)}}$ ENTSPRINGENDEN FUNCTIONEN.
\end{center}

\noindent\hspace*{\parindent}Zufolge einer Tagebuchnotiz hat \textsc{Gauss} die Inversion des hier behandelten
Integrals in die am Schlusse gewonnene Potenzreihe bereits am 9.\ September 1796 durchgeführt.

Der entwickelte Ansatz steht zur \textsc{Weierstrass}’schen Theorie in entsprechender
Beziehung, wie die Untersuchung des lemniscatischen Integrals zur \textsc{Jacobi}schen
Theorie. Liefert letzteres Integral den sogen. »harmonischen« Fall des elliptischen
Gebildes, so hat man es hier mit dem »aequianharmonischen« Falle zu thun.

Die \textsc{Gauss}’schen Angaben gehen sämmtlich leicht aus dem Additionstheorem
für $\boldsymbol{P}(\varphi)$ hervor, d.\,i. aus der oben für
$\boldsymbol{P}(3\psi+\varphi+\varphi')$ gegebenen Formel, welche, etwas anders geschrieben, so lautet:
\[
\boldsymbol{P}(3\psi+\varphi+\varphi')=
\frac{\begin{gathered}\boldsymbol{P}(\varphi)\boldsymbol{P}(\varphi')
(\boldsymbol{P}(\varphi)+\boldsymbol{P}(\varphi'))\\{}
+2-2\sqrt{1+\boldsymbol{P}(\varphi)^3}\sqrt{1+\boldsymbol{P}(\varphi')^3}\end{gathered}}
{(\boldsymbol{P}(\varphi)-\boldsymbol{P}(\varphi'))^2},
\]
\noindent und welche \textsc{Gauss} zweifellos auf dem auch beim Lemniscatenintegral befolgten
Wege, d.\,i. vermöge des ursprünglichen \textsc{Euler}schen Gedankenganges gewonnen hat.
Übrigens kommt diese Formel unmittelbar auf das Additionstheorem der
\textsc{Weierstrass}’schen Function $\wp(u)$ für $g_2=0$, $g_3=-4$ zurück.

\Needspace{12\baselineskip}
Aus dem Additionstheorem ergeben sich leicht die Formeln:
\[
\boldsymbol{P}(2\varphi+3\psi)=\frac{\boldsymbol{P}(\varphi)(\boldsymbol{P}(\varphi)^3-8)}
{4(\boldsymbol{P}(\varphi)^3+1)},
\]
\[
\boldsymbol{P}(3\varphi)=\frac{\boldsymbol{P}(\varphi)^9-96\boldsymbol{P}(\varphi)^6+48\boldsymbol{P}(\varphi)^3+64}
{9\boldsymbol{P}(\varphi)^2(\boldsymbol{P}(\varphi)^3+4)^2}.
\]
\noindent Mit Hülfe dieser Gleichungen entwickelt man von $\boldsymbol{P}(0)=-1$,
$\boldsymbol{P}(\psi)=0$ aus ohne Mühe die weiteren \textsc{Gauss}’schen Angaben.\hfill\textsc{Fricke.}
